\section{RAG-Systemarchitektur
(Event-Driven)}\label{rag-systemarchitektur-event-driven}

\subsection{Gliederung}\label{gliederung}

\begin{itemize}
\tightlist
\item
  \hyperref[1-zielsetzung]{1. Zielsetzung}

  \begin{itemize}
  \tightlist
  \item
    \hyperref[11-zielgruppe--lesehinweise]{1.1 Zielgruppe \&
    Lesehinweise}
  \end{itemize}
\item
  \hyperref[2-service-landschaft]{2. Service-Landschaft}

  \begin{itemize}
  \tightlist
  \item
    \hyperref[21-externe-abhuxe4ngigkeiten]{2.1 Externe Abhängigkeiten}
  \item
    \hyperref[22-schnittstellen--kontrakte-uxfcberblick]{2.2
    Schnittstellen \& Kontrakte (Überblick)}
  \end{itemize}
\item
  \hyperref[3-architektur-diagramme]{3. Architektur-Diagramme}

  \begin{itemize}
  \tightlist
  \item
    \hyperref[31-ingestion--dokumenten-pipeline-mit-kafka]{3.1 Ingestion
    / Dokumenten-Pipeline (mit Kafka)}
  \item
    \hyperref[32-inference--query-pipeline-ohne-kafka]{3.2 Inference /
    Query-Pipeline (ohne Kafka)}
  \item
    \hyperref[33-gesamtuxfcbersicht-end-to-end-flow-mit-topics]{3.3
    Gesamtübersicht: End-to-End Flow mit Topics}
  \item
    \hyperref[34-sequenzdiagramm-ingestion--retrieval-synchron--asynchron]{3.4
    Sequenzdiagramm: Ingestion \& Retrieval}
  \item
    \hyperref[35-visualisierung-von-topics--messages]{3.5 Visualisierung
    von Topics \& Messages}
  \item
    \hyperref[354-topic-vs-messages--abgegrenzt-szenario-document-received]{3.5.4
    Topic vs.~Messages -- abgegrenzt}
  \end{itemize}
\item
  \hyperref[4-kafka-nervensystem-topic-strategie]{4. Kafka
  ``Nervensystem'' (Topic-Strategie)}

  \begin{itemize}
  \tightlist
  \item
    \hyperref[41-topic-uxfcbersicht]{4.1 Topic-Übersicht}
  \item
    \hyperref[42-topic-payloads-schnittstellen-kontrakt]{4.2
    Topic-Payloads (Schnittstellen-Kontrakt)}
  \item
    \hyperref[43-beispiel-events-je-topic]{4.3 Beispiel-Events je Topic}
  \item
    \hyperref[44-producer-topic-consumer-mapping]{4.4
    Producer--Topic--Consumer Mapping}
  \item
    \hyperref[45-architekten-notiz-json-vs-avro-bewusst-json]{4.5
    Architekten-Notiz: JSON vs.~Avro (bewusst JSON)}
  \end{itemize}
\item
  \hyperref[5-end-to-end-datenfluss-pipelines]{5. End-to-End Datenfluss
  (Pipelines)}

  \begin{itemize}
  \tightlist
  \item
    \hyperref[51-ingestion-pipeline-daten-aufnahme]{5.1 Ingestion
    Pipeline (Daten-Aufnahme)}
  \item
    \hyperref[52-inference-pipeline-abfragechat]{5.2 Inference Pipeline
    (Abfrage/Chat)}
  \item
    \hyperref[53-wichtige-rest-endpunkte-schnittstellen-kontrakt]{5.3
    Wichtige REST-Endpunkte}
  \item
    \hyperref[54-beispielszenario-ingestion--retrieval-end-to-end]{5.4
    Beispielszenario: Ingestion \& Retrieval End-to-End}
  \item
    \hyperref[55-mandantentrennung--deduplizierung-tenantid--contenthash]{5.5
    Mandantentrennung \& Deduplizierung (\texttt{tenantId} \&
    \texttt{contentHash})}
  \end{itemize}
\item
  \hyperref[6-warum-dieses-konzept-argumentation]{6. Warum dieses
  Konzept? (Argumentation)}
\item
  \hyperref[7-betrieb-skalierung--bottlenecks]{7. Betrieb, Skalierung \&
  Bottlenecks}

  \begin{itemize}
  \tightlist
  \item
    \hyperref[71-wie-wir-topics-verstehen]{7.1 Wie wir Topics verstehen}
  \item
    \hyperref[72-ablauf-bei-multi-upload-viele-dokumente-auf-einmal]{7.2
    Ablauf bei Multi-Upload}
  \item
    \hyperref[73-retry-mechanismen--fehlerbehandlung]{7.3
    Retry-Mechanismen \& Fehlerbehandlung}
  \item
    \hyperref[74-typische-bottlenecks-pro-stage]{7.4 Typische
    Bottlenecks pro Stage}
  \item
    \hyperref[75-zusammenspiel-mit-kubernetes]{7.5 Zusammenspiel mit
    Kubernetes}
  \item
    \hyperref[76-fehlerszenarien--reprocessing-strategien]{7.6
    Fehlerszenarien \& Reprocessing-Strategien}
  \end{itemize}
\item
  \hyperref[8-glossar-wichtige-begriffe]{8. Glossar (wichtige Begriffe)}
\item
  \hyperref[9-weiterfuxfchrende-literatur--referenzen]{9. Weiterführende
  Literatur \& Referenzen}
\end{itemize}

\begin{center}\rule{0.5\linewidth}{0.5pt}\end{center}

\subsection{1. Zielsetzung}\label{zielsetzung}

Dieses Dokument beschreibt die \textbf{Architektur unserer RAG-Pipeline}
als event-getriebenes, mandantenfähiges System:

\begin{itemize}
\tightlist
\item
  \textbf{Dokumenten-Ingestion} (Extraktion, Anreicherung, Embeddings)
\item
  \textbf{Wissenssuche} (Vektorsuche \& Reranking)
\item
  \textbf{Antwortgenerierung} (LLM)
\end{itemize}

Im Fokus stehen \textbf{Services, Topics, Datenflüsse und externe
Abhängigkeiten} -- also wie das Projekt logisch aufgebaut ist.

\subsection{1.1 Zielgruppe \&
Lesehinweise}\label{zielgruppe-lesehinweise}

\begin{itemize}
\tightlist
\item
  \textbf{Für wen}:

  \begin{itemize}
  \tightlist
  \item
    \textbf{Backend-/Data-Engineers}: Architektur, Services, Payloads
    (Abschnitte 2--5).
  \item
    \textbf{SRE/DevOps}: Betrieb, Skalierung, Bottlenecks (Abschnitt 7).
  \item
    \textbf{Produkt/Stakeholder}: High-Level-Überblick, „Warum dieses
    Konzept?{\kern0pt}`` (Abschnitt 6).
  \end{itemize}
\item
  \textbf{Wie lesen}:

  \begin{itemize}
  \tightlist
  \item
    Wenn du \textbf{neu im Projekt} bist, starte mit Abschnitt 2 und 3.
  \item
    Wenn du \textbf{einen Incident} oder Bottleneck analysierst, springe
    direkt zu Abschnitt 7.
  \item
    Wenn du \textbf{Argumente für Architekturentscheidungen} brauchst,
    lies Abschnitt 6.
  \end{itemize}
\end{itemize}

\begin{center}\rule{0.5\linewidth}{0.5pt}\end{center}

\subsection{2. Service-Landschaft}\label{service-landschaft}

{\def\LTcaptype{none} % do not increment counter
\begin{longtable}[]{@{}
  >{\raggedright\arraybackslash}p{(\linewidth - 4\tabcolsep) * \real{0.2340}}
  >{\raggedright\arraybackslash}p{(\linewidth - 4\tabcolsep) * \real{0.1170}}
  >{\raggedright\arraybackslash}p{(\linewidth - 4\tabcolsep) * \real{0.6489}}@{}}
\toprule\noalign{}
\begin{minipage}[b]{\linewidth}\raggedright
Service
\end{minipage} & \begin{minipage}[b]{\linewidth}\raggedright
Sprache
\end{minipage} & \begin{minipage}[b]{\linewidth}\raggedright
Hauptaufgabe
\end{minipage} \\
\midrule\noalign{}
\endhead
\bottomrule\noalign{}
\endlastfoot
\textbf{Middleware} & Spring Boot & API Gateway, Authentifizierung,
S3-Management, Proxy-Logik. \\
\textbf{Extraction Service} & Python & Text-Extraktion aus PDFs, Bildern
und Office-Dokumenten. \\
\textbf{Miner Service} & Python & Daten-Mining: Extraktion von Entitäten
(Daten, IDs, Beträge). \\
\textbf{AI Service} & Python & Text-Chunking und Erstellung von
Embeddings via Azure OpenAI. \\
\textbf{Data Service} & Python & Orchestrierung der Datenflüsse und
Zustandsüberwachung. \\
\end{longtable}
}

\subsubsection{2.1 Externe
Abhängigkeiten}\label{externe-abhuxe4ngigkeiten}

{\def\LTcaptype{none} % do not increment counter
\begin{longtable}[]{@{}
  >{\raggedright\arraybackslash}p{(\linewidth - 4\tabcolsep) * \real{0.1449}}
  >{\raggedright\arraybackslash}p{(\linewidth - 4\tabcolsep) * \real{0.2174}}
  >{\raggedright\arraybackslash}p{(\linewidth - 4\tabcolsep) * \real{0.6377}}@{}}
\toprule\noalign{}
\begin{minipage}[b]{\linewidth}\raggedright
Service
\end{minipage} & \begin{minipage}[b]{\linewidth}\raggedright
Externer Dienst
\end{minipage} & \begin{minipage}[b]{\linewidth}\raggedright
Zweck
\end{minipage} \\
\midrule\noalign{}
\endhead
\bottomrule\noalign{}
\endlastfoot
Middleware & S3 & Persistenz der Original-Dokumente \\
Middleware & Azure OpenAI & Embeddings \& LLM-Inferenz (Abfragepfad) \\
AI Service & Azure OpenAI & Embeddings (Ingestionpfad) \\
AI Service & Milvus & Schreiben der Vektoren \\
Middleware & Milvus & Lesen der Vektoren für die Ähnlichkeitssuche \\
\end{longtable}
}

\subsubsection{2.2 Schnittstellen \& Kontrakte
(Überblick)}\label{schnittstellen-kontrakte-uxfcberblick}

\begin{itemize}
\tightlist
\item
  \textbf{REST-API (Middleware)}: Upload von Dokumenten, Starten von
  Pipelines, Status-Abfragen, Chat-/Query-Endpunkte.
\item
  \textbf{Kafka-Topics}: Asynchrone Übergabe zwischen Services; jedes
  Topic hat einen klar definierten Payload-Kontrakt.
\item
  \textbf{Data Service}: Zentrale Fassade für Milvus (Schreiben/Lesen
  von Vektoren) und Status-Updates in Postgres.
\item
  \textbf{AI Service}: Spricht ausschließlich mit Azure OpenAI und dem
  Data Service -- nie direkt mit Milvus.
\end{itemize}

\begin{center}\rule{0.5\linewidth}{0.5pt}\end{center}

\subsection{3. Architektur-Diagramme}\label{architektur-diagramme}

Um den Ablauf klar zu trennen, zeigen wir zwei Sichten:

\begin{itemize}
\tightlist
\item
  eine \textbf{Ingestion-Sicht} mit Kafka-Topics und Python-Services und
\item
  eine \textbf{Inference-Sicht} für Queries über Middleware, Azure
  OpenAI und Milvus.
\end{itemize}

\subsubsection{3.1 Ingestion / Dokumenten-Pipeline (mit
Kafka)}\label{ingestion-dokumenten-pipeline-mit-kafka}

\begin{figure}
\centering
\pandocbounded{\includesvg[keepaspectratio]{diagrams/diagram-01.svg}}
\caption{Diagramm 1}
\end{figure}

\subsubsection{3.2 Inference / Query-Pipeline (ohne
Kafka)}\label{inference-query-pipeline-ohne-kafka}

\begin{figure}
\centering
\pandocbounded{\includesvg[keepaspectratio]{diagrams/diagram-02.svg}}
\caption{Diagramm 2}
\end{figure}

\subsubsection{3.3 Gesamtübersicht: End-to-End Flow mit
Topics}\label{gesamtuxfcbersicht-end-to-end-flow-mit-topics}

Die runden Knoten unten repräsentieren \textbf{Kafka-Topics}.\\
In den Labels ist jeweils kurz angedeutet, \textbf{welche Payload} darin
steckt.

\begin{figure}
\centering
\pandocbounded{\includesvg[keepaspectratio]{diagrams/diagram-03.svg}}
\caption{Diagramm 3}
\end{figure}

\begin{center}\rule{0.5\linewidth}{0.5pt}\end{center}

\subsubsection{3.4 Sequenzdiagramm: Ingestion \& Retrieval (synchron +
asynchron)}\label{sequenzdiagramm-ingestion-retrieval-synchron-asynchron}

Das folgende Sequenzdiagramm stellt den Gesamtfluss noch einmal aus
Sicht der \textbf{Aufrufe und Events} dar.

\begin{figure}
\centering
\pandocbounded{\includesvg[keepaspectratio]{diagrams/diagram-04.svg}}
\caption{Diagramm 4}
\end{figure}

\subsubsection{3.5 Visualisierung von Topics \&
Messages}\label{visualisierung-von-topics-messages}

\paragraph{3.5.1 Ein Topic als Log
(Zeitachse)}\label{ein-topic-als-log-zeitachse}

\begin{figure}
\centering
\pandocbounded{\includesvg[keepaspectratio]{diagrams/diagram-05.svg}}
\caption{Diagramm 5}
\end{figure}

\paragraph{3.5.2 Ein Dokument ueber mehrere Topics
hinweg}\label{ein-dokument-ueber-mehrere-topics-hinweg}

\begin{figure}
\centering
\pandocbounded{\includesvg[keepaspectratio]{diagrams/diagram-06.svg}}
\caption{Diagramm 6}
\end{figure}

\paragraph{3.5.3 Producer--Topic--Consumer auf einen
Blick}\label{producertopicconsumer-auf-einen-blick}

\begin{figure}
\centering
\pandocbounded{\includesvg[keepaspectratio]{diagrams/diagram-07.svg}}
\caption{Diagramm 7}
\end{figure}

\paragraph{\texorpdfstring{3.5.4 Topic vs.~Messages -- abgegrenzt
(Szenario
\texttt{document-received})}{3.5.4 Topic vs.~Messages -- abgegrenzt (Szenario document-received)}}\label{topic-vs.-messages-abgegrenzt-szenario-document-received}

\textbf{Was ist was?} Ein \textbf{Topic} ist ein benannter, append-only
\textbf{Log}; eine \textbf{Message} ist ein einzelnes Event (ein
JSON-Payload) an einer \textbf{Offset}-Position. Jede Message gehört zu
genau einem Topic und wird von Producer(s) geschrieben und von
Consumer(s) gelesen.

\begin{figure}
\centering
\pandocbounded{\includesvg[keepaspectratio]{diagrams/diagram-08.svg}}
\caption{Diagramm 8}
\end{figure}

\textbf{Abgrenzung im Überblick}

{\def\LTcaptype{none} % do not increment counter
\begin{longtable}[]{@{}
  >{\raggedright\arraybackslash}p{(\linewidth - 2\tabcolsep) * \real{0.2051}}
  >{\raggedright\arraybackslash}p{(\linewidth - 2\tabcolsep) * \real{0.7949}}@{}}
\toprule\noalign{}
\begin{minipage}[b]{\linewidth}\raggedright
Begriff
\end{minipage} & \begin{minipage}[b]{\linewidth}\raggedright
Bedeutung in unserem Szenario
\end{minipage} \\
\midrule\noalign{}
\endhead
\bottomrule\noalign{}
\endlastfoot
\textbf{Topic} & Ein \textbf{Kanal} (z.\,B. \texttt{document-received}).
Enthält alle Messages in Reihenfolge; wird nicht gelöscht, nur angehängt
(bis Retention). \\
\textbf{Message} & Ein \textbf{Event} = ein JSON-Objekt (Payload) mit
z.\,B. \texttt{documentId}, \texttt{tenantId}, \texttt{s3Path},
\texttt{contentHash}. Eine Message = ein hochgeladenes Dokument, das zur
Extraktion ansteht. \\
\textbf{Offset} & Eindeutige \textbf{Position} der Message im Topic-Log.
Consumer merken sich den letzten gelesenen Offset (commit). \\
\textbf{Partition} (optional) & Bei Skalierung kann ein Topic in
Partitionen unterteilt werden; z.\,B. Partition-Key = \texttt{tenantId}
für Mandantentrennung und paralleles Lesen. \\
\end{longtable}
}

\subsection{4. Kafka ``Nervensystem''
(Topic-Strategie)}\label{kafka-nervensystem-topic-strategie}

Die Kommunikation erfolgt über spezialisierte Topics. Jede Nachricht
trägt im Header oder Payload die \texttt{tenant\_id} zur Sicherstellung
der Daten-Isolation.

\subsubsection{4.1 Topic-Übersicht}\label{topic-uxfcbersicht}

\begin{enumerate}
\def\labelenumi{\arabic{enumi}.}
\tightlist
\item
  \textbf{Topic:} \texttt{document-received}
\end{enumerate}

\begin{itemize}
\tightlist
\item
  \textbf{Trigger:} Spring Boot empfängt Datei.
\item
  \textbf{Inhalt:} S3-Link zum Original, Metadaten, Mandanten-ID.
\end{itemize}

\begin{enumerate}
\def\labelenumi{\arabic{enumi}.}
\setcounter{enumi}{1}
\tightlist
\item
  \textbf{Topic:} \texttt{document-to-extract}
\end{enumerate}

\begin{itemize}
\tightlist
\item
  \textbf{Trigger:} Data Service gibt Dokument zur Verarbeitung frei.
\item
  \textbf{Inhalt:} Pfad zur Rohdatei.
\end{itemize}

\begin{enumerate}
\def\labelenumi{\arabic{enumi}.}
\setcounter{enumi}{2}
\tightlist
\item
  \textbf{Topic:} \texttt{content-extracted}
\end{enumerate}

\begin{itemize}
\tightlist
\item
  \textbf{Trigger:} Extraction Service fertig.
\item
  \textbf{Inhalt:} extrahierter Reintext (Raw-Text).
\end{itemize}

\begin{enumerate}
\def\labelenumi{\arabic{enumi}.}
\setcounter{enumi}{3}
\tightlist
\item
  \textbf{Topic:} \texttt{metadata-enriched}
\end{enumerate}

\begin{itemize}
\tightlist
\item
  \textbf{Trigger:} Miner Service hat Fachdaten (z.\,B. Rechnungsdatum)
  gefunden.
\item
  \textbf{Inhalt:} Text + strukturierte Zusatzdaten.
\end{itemize}

\begin{enumerate}
\def\labelenumi{\arabic{enumi}.}
\setcounter{enumi}{4}
\tightlist
\item
  \textbf{Topic:} \texttt{vector-ready-to-index}
\end{enumerate}

\begin{itemize}
\tightlist
\item
  \textbf{Trigger:} AI Service hat Text gechunked und in Vektoren
  umgewandelt.
\item
  \textbf{Inhalt:} Vektor-Arrays, Chunks, finale Metadaten \&
  \texttt{tenant\_id}.
\end{itemize}

\subsubsection{4.2 Topic-Payloads
(Schnittstellen-Kontrakt)}\label{topic-payloads-schnittstellen-kontrakt}

Alle Payloads sind JSON-Objekte; Pflichtfelder sind fett markiert.

\begin{itemize}
\tightlist
\item
  \textbf{Topic \texttt{document-received}}

  \begin{itemize}
  \tightlist
  \item
    \textbf{\texttt{documentId}}: eindeutige ID des Dokuments
  \item
    \textbf{\texttt{tenantId}}: Mandanten-ID
  \item
    \textbf{\texttt{s3Path}}: Pfad zur Originaldatei in S3
  \item
    \texttt{batchId}: optionale Gruppen-ID für Bulk-Uploads
  \item
    \texttt{filename}, \texttt{mimeType}, \texttt{uploadedBy}
  \item
    \texttt{contentHash}: Hash des Datei-Inhalts (z.\,B. SHA-256)
  \end{itemize}
\item
  \textbf{Topic \texttt{document-to-extract}}

  \begin{itemize}
  \tightlist
  \item
    \textbf{\texttt{documentId}}, \textbf{\texttt{tenantId}}
  \item
    \textbf{\texttt{s3Path}}
  \item
    \texttt{priority}: z.\,B. „LOW/MEDIUM/HIGH``
  \end{itemize}
\item
  \textbf{Topic \texttt{content-extracted}}

  \begin{itemize}
  \tightlist
  \item
    \textbf{\texttt{documentId}}, \textbf{\texttt{tenantId}}
  \item
    \textbf{\texttt{plainText}}: extrahierter Volltext
  \item
    \texttt{structure}: optionale Layout-/Tabelleninfos (z.\,B. von
    Document Intelligence)
  \item
    \texttt{language}, \texttt{sourceFileName}
  \end{itemize}
\item
  \textbf{Topic \texttt{metadata-enriched}}

  \begin{itemize}
  \tightlist
  \item
    \textbf{\texttt{documentId}}, \textbf{\texttt{tenantId}}
  \item
    \textbf{\texttt{plainText}}
  \item
    \textbf{\texttt{metadata}}: strukturierte Fachdaten (z.\,B.
    Rechnungsnummer, Datum, Beträge)
  \item
    \texttt{detectedEntities}: optionale, modellbasierte
    Extraktionsergebnisse
  \end{itemize}
\item
  \textbf{Topic \texttt{vector-ready-to-index}}

  \begin{itemize}
  \tightlist
  \item
    \textbf{\texttt{documentId}}, \textbf{\texttt{tenantId}}
  \item
    \textbf{\texttt{chunks}}: Liste von Chunks mit:

    \begin{itemize}
    \tightlist
    \item
      \textbf{\texttt{chunkId}}, \textbf{\texttt{text}}
    \item
      \textbf{\texttt{vector}}: Embedding-Vektor
    \item
      \texttt{page}, \texttt{offset}, \texttt{section},
      \texttt{metadata}
    \end{itemize}
  \item
    \texttt{indexName}: logischer Index-/Collection-Name in Milvus
  \end{itemize}
\end{itemize}

\subsubsection{4.3 Beispiel-Events je
Topic}\label{beispiel-events-je-topic}

\textbf{\texttt{document-received}}

\begin{Shaded}
\begin{Highlighting}[]
\FunctionTok{\{}
  \DataTypeTok{"documentId"}\FunctionTok{:} \StringTok{"DOC{-}123"}\FunctionTok{,}
  \DataTypeTok{"tenantId"}\FunctionTok{:} \StringTok{"TENANT{-}A"}\FunctionTok{,}
  \DataTypeTok{"batchId"}\FunctionTok{:} \StringTok{"BATCH{-}2026{-}03{-}01{-}001"}\FunctionTok{,}
  \DataTypeTok{"s3Path"}\FunctionTok{:} \StringTok{"s3://bucket/raw/TENANT{-}A/DOC{-}123.pdf"}\FunctionTok{,}
  \DataTypeTok{"filename"}\FunctionTok{:} \StringTok{"rechnung\_123.pdf"}\FunctionTok{,}
  \DataTypeTok{"mimeType"}\FunctionTok{:} \StringTok{"application/pdf"}\FunctionTok{,}
  \DataTypeTok{"uploadedBy"}\FunctionTok{:} \StringTok{"user@kunde.de"}\FunctionTok{,}
  \DataTypeTok{"contentHash"}\FunctionTok{:} \StringTok{"b8a6b4c0...ef"}
\FunctionTok{\}}
\end{Highlighting}
\end{Shaded}

\textbf{\texttt{document-to-extract}}

\begin{Shaded}
\begin{Highlighting}[]
\FunctionTok{\{}
  \DataTypeTok{"documentId"}\FunctionTok{:} \StringTok{"DOC{-}123"}\FunctionTok{,}
  \DataTypeTok{"tenantId"}\FunctionTok{:} \StringTok{"TENANT{-}A"}\FunctionTok{,}
  \DataTypeTok{"s3Path"}\FunctionTok{:} \StringTok{"s3://bucket/raw/TENANT{-}A/DOC{-}123.pdf"}\FunctionTok{,}
  \DataTypeTok{"priority"}\FunctionTok{:} \StringTok{"MEDIUM"}
\FunctionTok{\}}
\end{Highlighting}
\end{Shaded}

\textbf{\texttt{content-extracted}}

\begin{Shaded}
\begin{Highlighting}[]
\FunctionTok{\{}
  \DataTypeTok{"documentId"}\FunctionTok{:} \StringTok{"DOC{-}123"}\FunctionTok{,}
  \DataTypeTok{"tenantId"}\FunctionTok{:} \StringTok{"TENANT{-}A"}\FunctionTok{,}
  \DataTypeTok{"plainText"}\FunctionTok{:} \StringTok{"Rechnung 123}\CharTok{\textbackslash{}n}\StringTok{Kunde XY GmbH}\CharTok{\textbackslash{}n}\StringTok{Betrag: 1.234,56 EUR}\CharTok{\textbackslash{}n}\StringTok{..."}\FunctionTok{,}
  \DataTypeTok{"structure"}\FunctionTok{:} \FunctionTok{\{}
    \DataTypeTok{"pages"}\FunctionTok{:} \DecValTok{3}\FunctionTok{,}
    \DataTypeTok{"hasTables"}\FunctionTok{:} \KeywordTok{true}
  \FunctionTok{\},}
  \DataTypeTok{"language"}\FunctionTok{:} \StringTok{"de"}\FunctionTok{,}
  \DataTypeTok{"sourceFileName"}\FunctionTok{:} \StringTok{"rechnung\_123.pdf"}
\FunctionTok{\}}
\end{Highlighting}
\end{Shaded}

\textbf{\texttt{metadata-enriched}}

\begin{Shaded}
\begin{Highlighting}[]
\FunctionTok{\{}
  \DataTypeTok{"documentId"}\FunctionTok{:} \StringTok{"DOC{-}123"}\FunctionTok{,}
  \DataTypeTok{"tenantId"}\FunctionTok{:} \StringTok{"TENANT{-}A"}\FunctionTok{,}
  \DataTypeTok{"plainText"}\FunctionTok{:} \StringTok{"Rechnung 123}\CharTok{\textbackslash{}n}\StringTok{Kunde XY GmbH}\CharTok{\textbackslash{}n}\StringTok{Betrag: 1.234,56 EUR}\CharTok{\textbackslash{}n}\StringTok{..."}\FunctionTok{,}
  \DataTypeTok{"metadata"}\FunctionTok{:} \FunctionTok{\{}
    \DataTypeTok{"invoiceNumber"}\FunctionTok{:} \StringTok{"123"}\FunctionTok{,}
    \DataTypeTok{"invoiceDate"}\FunctionTok{:} \StringTok{"2026{-}02{-}28"}\FunctionTok{,}
    \DataTypeTok{"amount"}\FunctionTok{:} \FloatTok{1234.56}\FunctionTok{,}
    \DataTypeTok{"currency"}\FunctionTok{:} \StringTok{"EUR"}\FunctionTok{,}
    \DataTypeTok{"customerName"}\FunctionTok{:} \StringTok{"XY GmbH"}
  \FunctionTok{\},}
  \DataTypeTok{"detectedEntities"}\FunctionTok{:} \OtherTok{[}
    \FunctionTok{\{} \DataTypeTok{"type"}\FunctionTok{:} \StringTok{"INVOICE\_NUMBER"}\FunctionTok{,} \DataTypeTok{"value"}\FunctionTok{:} \StringTok{"123"} \FunctionTok{\}}\OtherTok{,}
    \FunctionTok{\{} \DataTypeTok{"type"}\FunctionTok{:} \StringTok{"AMOUNT"}\FunctionTok{,} \DataTypeTok{"value"}\FunctionTok{:} \StringTok{"1234.56"} \FunctionTok{\}}
  \OtherTok{]}
\FunctionTok{\}}
\end{Highlighting}
\end{Shaded}

\textbf{\texttt{vector-ready-to-index}}

\begin{Shaded}
\begin{Highlighting}[]
\FunctionTok{\{}
  \DataTypeTok{"documentId"}\FunctionTok{:} \StringTok{"DOC{-}123"}\FunctionTok{,}
  \DataTypeTok{"tenantId"}\FunctionTok{:} \StringTok{"TENANT{-}A"}\FunctionTok{,}
  \DataTypeTok{"indexName"}\FunctionTok{:} \StringTok{"tenant{-}a{-}main"}\FunctionTok{,}
  \DataTypeTok{"chunks"}\FunctionTok{:} \OtherTok{[}
    \FunctionTok{\{}
      \DataTypeTok{"chunkId"}\FunctionTok{:} \StringTok{"DOC{-}123{-}0001"}\FunctionTok{,}
      \DataTypeTok{"text"}\FunctionTok{:} \StringTok{"Diese Rechnung bezieht sich auf den Vertrag 9876 ..."}\FunctionTok{,}
      \DataTypeTok{"vector"}\FunctionTok{:} \OtherTok{[}\FloatTok{0.0123}\OtherTok{,} \FloatTok{{-}0.9987}\OtherTok{,} \FloatTok{0.4532}\OtherTok{]}\FunctionTok{,}
      \DataTypeTok{"page"}\FunctionTok{:} \DecValTok{1}\FunctionTok{,}
      \DataTypeTok{"offset"}\FunctionTok{:} \DecValTok{0}\FunctionTok{,}
      \DataTypeTok{"section"}\FunctionTok{:} \StringTok{"header"}\FunctionTok{,}
      \DataTypeTok{"metadata"}\FunctionTok{:} \FunctionTok{\{}
        \DataTypeTok{"invoiceNumber"}\FunctionTok{:} \StringTok{"123"}\FunctionTok{,}
        \DataTypeTok{"invoiceDate"}\FunctionTok{:} \StringTok{"2026{-}02{-}28"}
      \FunctionTok{\}}
    \FunctionTok{\}}
  \OtherTok{]}
\FunctionTok{\}}
\end{Highlighting}
\end{Shaded}

\subsubsection{4.4 Producer--Topic--Consumer
Mapping}\label{producertopicconsumer-mapping}

{\def\LTcaptype{none} % do not increment counter
\begin{longtable}[]{@{}
  >{\raggedright\arraybackslash}p{(\linewidth - 6\tabcolsep) * \real{0.1314}}
  >{\raggedright\arraybackslash}p{(\linewidth - 6\tabcolsep) * \real{0.1679}}
  >{\raggedright\arraybackslash}p{(\linewidth - 6\tabcolsep) * \real{0.3723}}
  >{\raggedright\arraybackslash}p{(\linewidth - 6\tabcolsep) * \real{0.3285}}@{}}
\toprule\noalign{}
\begin{minipage}[b]{\linewidth}\raggedright
Producer
\end{minipage} & \begin{minipage}[b]{\linewidth}\raggedright
Topic
\end{minipage} & \begin{minipage}[b]{\linewidth}\raggedright
Consumer
\end{minipage} & \begin{minipage}[b]{\linewidth}\raggedright
Zweck
\end{minipage} \\
\midrule\noalign{}
\endhead
\bottomrule\noalign{}
\endlastfoot
Middleware & \texttt{document-received} & Data Service & Neues Dokument
ist hochgeladen \\
Data Service & \texttt{document-to-extract} & Extraction Service &
Dokument soll in Text extrahiert werden \\
Extraction Service & \texttt{content-extracted} & Miner Service &
Rohtext liegt zur fachlichen Anreicherung vor \\
Miner Service & \texttt{metadata-enriched} & AI Service & Angereicherter
Text ist bereit für Embedding \\
AI Service & \texttt{vector-ready-to-index} & (Indexer / Milvus-Writer,
Teil des AI/Data Service) & Vektoren können persistiert werden \\
\end{longtable}
}

\subsubsection{4.5 Architekten-Notiz: JSON vs.~Avro (bewusst
JSON)}\label{architekten-notiz-json-vs.-avro-bewusst-json}

\textbf{Kurz-Zusammenfassung:}\\
Alle Kafka-Payloads nutzen \textbf{JSON} -- keine Avro-Schemas. Das ist
eine bewusste Architekturentscheidung: Lesbarkeit, Debugging und
Interoperabilität wiegen derzeit höher als die theoretischen
Performance-Vorteile von Avro.

{\def\LTcaptype{none} % do not increment counter
\begin{longtable}[]{@{}
  >{\raggedright\arraybackslash}p{(\linewidth - 4\tabcolsep) * \real{0.2340}}
  >{\raggedright\arraybackslash}p{(\linewidth - 4\tabcolsep) * \real{0.3404}}
  >{\raggedright\arraybackslash}p{(\linewidth - 4\tabcolsep) * \real{0.4255}}@{}}
\toprule\noalign{}
\begin{minipage}[b]{\linewidth}\raggedright
Kriterium
\end{minipage} & \begin{minipage}[b]{\linewidth}\raggedright
JSON (gewählt)
\end{minipage} & \begin{minipage}[b]{\linewidth}\raggedright
Avro (Alternative)
\end{minipage} \\
\midrule\noalign{}
\endhead
\bottomrule\noalign{}
\endlastfoot
\textbf{Lesbarkeit} & Menschenlesbar in Kafka-UI, Logs, Admin-Tools &
Binär; Schema/Codegen nötig zum Inspizieren \\
\textbf{Debugging \& Onboarding} & Events sofort verständlich; geringer
Einstieg für alle Personas & Höhere Hürde für SRE/Fachseite ohne
Avro-Tooling \\
\textbf{Interoperabilität} & Universell (REST, Scripte, Metriken); kein
Schema Registry nötig & Schema Registry, Codegen, Build-Integration
erforderlich \\
\textbf{Performance / Effizienz} & Größere Messages, langsameres Parsing
bei sehr hohem Durchsatz & Kompakter, schnelleres Parsing -- relevant
bei Kafka als Engpass \\
\textbf{Kontrakt-Sicherheit} & Schemas in Doku (Abschnitt 4.2); weiche
Evolution möglich & Strikte Schemas, zentrale Evolution über Schema
Registry \\
\end{longtable}
}

\textbf{Warum JSON für diese Pipeline passt}

\begin{itemize}
\tightlist
\item
  Die \textbf{Bottlenecks} liegen typischerweise bei \textbf{Azure
  OpenAI}, \textbf{Document Intelligence} und \textbf{Milvus}, nicht
  beim JSON-Parsing auf dem Bus.
\item
  \textbf{Viele Personas} (Backend, Data, SRE, Fach) müssen Events
  schnell einordnen können -- JSON reduziert Reibung.
\item
  \textbf{DLQ-Analyse und Reprocessing} profitieren von lesbaren
  Payloads ohne zusätzliches Tooling.
\end{itemize}

\textbf{Operational Note}

\begin{itemize}
\tightlist
\item
  Die \textbf{Payload-Kontrakte} (Abschnitt 4.2) werden wie
  \textbf{strikte Schemas} gepflegt; Änderungen sind dokumentiert und
  rückwärtskompatibel zu handhaben.
\item
  Falls Kafka bei stark wachsendem Volumen zum Engpass wird:
  \textbf{Avro + Schema Registry} gezielt für High-Volume-Topics (z.\,B.
  \texttt{vector-ready-to-index}) evaluieren; optional \textbf{hybrid}:
  JSON an Rändern (REST, Admin), Avro intern für ausgewählte Topics.
\end{itemize}

\begin{center}\rule{0.5\linewidth}{0.5pt}\end{center}

\subsection{5. End-to-End Datenfluss
(Pipelines)}\label{end-to-end-datenfluss-pipelines}

\subsubsection{5.1 Ingestion Pipeline
(Daten-Aufnahme)}\label{ingestion-pipeline-daten-aufnahme}

\begin{enumerate}
\def\labelenumi{\arabic{enumi}.}
\tightlist
\item
  \textbf{Client} lädt Dokument über den \textbf{Spring Boot Proxy}
  hoch.
\item
  \textbf{Spring Boot} speichert Datei in \textbf{S3} und schreibt Event
  in \texttt{document-received}.
\item
  \textbf{Extraction Service} zieht Text aus dem Dokument
  (\texttt{document-to-extract} → \texttt{content-extracted}).
\item
  \textbf{Miner Service} reichert den Text mit Fachmetadaten an
  (\texttt{metadata-enriched}).
\item
  \textbf{AI Service} erstellt Chunks und ruft \textbf{Azure OpenAI} für
  Embeddings auf.
\item
  Die fertigen Vektoren werden in \textbf{Milvus} indiziert (mit
  \texttt{tenant\_id} als Partitions-Key).
\end{enumerate}

\subsubsection{5.2 Inference Pipeline
(Abfrage/Chat)}\label{inference-pipeline-abfragechat}

\begin{enumerate}
\def\labelenumi{\arabic{enumi}.}
\tightlist
\item
  \textbf{User} stellt eine Frage via \textbf{Spring Boot}.
\item
  \textbf{Spring Boot} wandelt die Frage via \textbf{Azure} in einen
  Vektor um.
\item
  \textbf{Spring Boot} führt eine Ähnlichkeitssuche in \textbf{Milvus}
  aus (Filter: \texttt{tenant\_id\ ==\ user\_tenant}).
\item
  Gefundene Fakten werden als Kontext an das \textbf{LLM (Azure)}
  gesendet.
\item
  Die Antwort wird an den Client zurückgegeben.
\end{enumerate}

\subsubsection{5.3 Wichtige REST-Endpunkte
(Schnittstellen-Kontrakt)}\label{wichtige-rest-endpunkte-schnittstellen-kontrakt}

\textbf{Upload / Pipeline-Start}

\begin{itemize}
\tightlist
\item
  \textbf{\texttt{POST\ /api/v1/pipelines/documents}}

  \begin{itemize}
  \tightlist
  \item
    \textbf{Body (JSON oder multipart + JSON-Metadaten):}

    \begin{itemize}
    \tightlist
    \item
      \texttt{tenantId} (string, Pflicht)
    \item
      \texttt{mode} (string, z.\,B. \texttt{"ALL\_INCLUSIVE"})
    \item
      \texttt{documents} (Liste) mit:

      \begin{itemize}
      \tightlist
      \item
        \texttt{filename}, \texttt{contentType}
      \item
        Binärinhalt oder Referenz (z.\,B. S3-PreSigned-URL)
      \end{itemize}
    \end{itemize}
  \item
    \textbf{Response 202 (Accepted):}

    \begin{itemize}
    \tightlist
    \item
      \texttt{batchId} (string)
    \item
      \texttt{acceptedCount} (int)
    \item
      \texttt{duplicatesSkipped} (int) -- Anzahl Dokumente, deren Inhalt
      bereits bekannt war
    \end{itemize}
  \end{itemize}
\end{itemize}

\textbf{Pipeline-Status abfragen}

\begin{itemize}
\tightlist
\item
  \textbf{\texttt{GET\ /api/v1/pipelines/\{batchId\}/status}}

  \begin{itemize}
  \tightlist
  \item
    \textbf{Response 200 (JSON):}

    \begin{itemize}
    \tightlist
    \item
      \texttt{batchId}
    \item
      \texttt{tenantId}
    \item
      \texttt{overallStatus} (z.\,B. \texttt{PENDING}, \texttt{RUNNING},
      \texttt{COMPLETED}, \texttt{FAILED})
    \item
      \texttt{documents} (Liste mit pro Dokument: \texttt{documentId},
      \texttt{status}, \texttt{errors}, \texttt{duplicateOfDocumentId?})
    \end{itemize}
  \end{itemize}
\end{itemize}

\textbf{Frage stellen (Chat / Retrieval)}

\begin{itemize}
\tightlist
\item
  \textbf{\texttt{POST\ /api/v1/chat}}

  \begin{itemize}
  \tightlist
  \item
    \textbf{Body:}

    \begin{itemize}
    \tightlist
    \item
      \texttt{tenantId} (string, Pflicht)
    \item
      \texttt{question} (string)
    \item
      \texttt{batchId} oder Filterkriterien (optional, um auf bestimmte
      Dokumentmengen zu beschränken)
    \end{itemize}
  \item
    \textbf{Response 200:}

    \begin{itemize}
    \tightlist
    \item
      \texttt{answer} (string)
    \item
      \texttt{sources} (Liste von Referenzen auf Chunks/Dokumente)
    \end{itemize}
  \end{itemize}
\end{itemize}

\subsubsection{5.4 Beispielszenario: Ingestion \& Retrieval
End-to-End}\label{beispielszenario-ingestion-retrieval-end-to-end}

Dieses Szenario verbindet die Ingestion- und Inference-Pipeline zu einem
realistischen Ablauf aus Sicht eines Mandanten.

\begin{enumerate}
\def\labelenumi{\arabic{enumi}.}
\tightlist
\item
  \textbf{Bulk-Upload vom Client (10.000 Dokumente)}

  \begin{itemize}
  \tightlist
  \item
    Der Client ruft die Middleware per REST auf (z.\,B. „ALL INCLUSIVE
    PIPELINE``) und übergibt bis zu 10.000 Dokumente -- entweder als
    \textbf{einzelne Requests} oder als \textbf{Batch}.\\
  \item
    Für jedes Dokument wird intern eine eigene \texttt{document\_id}
    vergeben; optional wird zusätzlich eine gemeinsame
    \texttt{batch\_id} geführt.
  \item
    Die Middleware bzw. der Data Service berechnet für jedes Dokument
    einen \textbf{Inhalts-Hash} (z.\,B.
    \texttt{content\_hash\ =\ SHA-256(binary)}).
  \end{itemize}
\item
  \textbf{Persistenz in S3 und Postgres durch den Data Service}

  \begin{itemize}
  \tightlist
  \item
    Die Middleware fungiert als Proxy und leitet die Uploads an den
    \textbf{Data Service} weiter.\\
  \item
    Der Data Service prüft anhand von \texttt{content\_hash}, ob der
    \textbf{identische Inhalt} für denselben \texttt{tenantId} bereits
    existiert:

    \begin{itemize}
    \tightlist
    \item
      Falls \textbf{neu}: Dokument wird im \textbf{S3 Storage}
      gespeichert, in \textbf{Postgres} entsteht ein neuer Eintrag
      (\texttt{document\_id}, \texttt{tenant\_id},
      \texttt{content\_hash}, S3-Pfad, Status „RECEIVED``).\\
    \item
      Falls \textbf{Duplikat}: Es wird nur eine \textbf{weitere logische
      Referenz} (z.\,B. neue \texttt{document\_id} mit Verweis auf
      bestehendes physisches Dokument) angelegt; teure
      Verarbeitungsschritte können übersprungen oder wiederverwendet
      werden.
    \end{itemize}
  \item
    Parallel erzeugt der Data Service Events in Kafka (z.\,B.
    \texttt{document-received} / \texttt{document-to-extract}), wie in
    Abschnitt 4.1 beschrieben -- Duplikate können hier je nach Policy
    entweder \textbf{nicht} mehr in die Pipeline gegeben oder als
    „Shortcut`` markiert werden (z.\,B. nur Metadaten aktualisieren).
  \end{itemize}
\item
  \textbf{Status-Transparenz für den Mandanten}

  \begin{itemize}
  \tightlist
  \item
    Nach erfolgreichem Upload erhält der Client eine
    \textbf{Pipeline-/Job-ID} (z.\,B. \texttt{batch\_id}), unter der
    alle Dokumente des Uploads verwaltet werden.\\
  \item
    Über einen dedizierten Status-Endpoint der Middleware kann der
    Client den \textbf{aktuellen Fortschritt} abfragen (z.\,B. „20 \%
    extrahiert, 10 \% indexiert, 70 \% pending``).\\
  \item
    Die Statusinformationen stammen aus Postgres, das vom Data Service
    bei jedem Verarbeitungsschritt aktualisiert wird (z.\,B.
    \texttt{EXTRACTED}, \texttt{EMBEDDED}, \texttt{INDEXED}).
  \end{itemize}
\item
  \textbf{Extraktion mit Azure Document Intelligence}

  \begin{itemize}
  \tightlist
  \item
    Nachdem ein Dokument im S3 Storage liegt und der entsprechende
    Kafka-Event verarbeitet wurde, zieht der \textbf{Extraction Service}
    das Dokument aus S3.\\
  \item
    Der Extraction Service ruft \textbf{Azure Document Intelligence} auf
    und erhält strukturierten Volltext, Layoutinformationen und ggf.
    Tabellen.\\
  \item
    Nach erfolgreicher Extraktion wird der Status in Postgres
    aktualisiert (z.\,B. \texttt{EXTRACTED}), und ein Event
    \texttt{content-extracted} wird erzeugt.
  \end{itemize}
\item
  \textbf{Speicherung von Volltext \& Chunks}

  \begin{itemize}
  \tightlist
  \item
    Aus der Antwort von Azure werden \textbf{Volltext und Chunks}
    erzeugt.\\
  \item
    Der Volltext (Extraktionsergebnis) wird erneut in S3 gespeichert
    (z.\,B. als Markdown/JSON), referenziert über Postgres.\\
  \item
    Die Chunks selbst werden in \textbf{Milvus} zunächst als reine
    Text/Metadaten-Einträge oder in separaten Sammlungen vorbereitet, je
    nach Implementierung.
  \end{itemize}
\item
  \textbf{Erstellung der Embeddings durch den AI Service}

  \begin{itemize}
  \tightlist
  \item
    Der \textbf{AI Service} übernimmt die vorbereiteten Chunks (über
    Kafka-Events wie \texttt{metadata-enriched} oder interne Abfragen)
    und sendet sie an \textbf{Azure OpenAI}, um Embeddings zu
    erzeugen.\\
  \item
    Der AI Service sorgt für Rate-Limiting und Batch-Verarbeitung, wie
    in Abschnitt 7 beschrieben.
  \end{itemize}
\item
  \textbf{Speicherung der Embeddings über den Data Service}

  \begin{itemize}
  \tightlist
  \item
    Der AI Service leitet die Embedding-Ergebnisse an den \textbf{Data
    Service} weiter oder schreibt sie direkt in Milvus und meldet den
    Erfolg an den Data Service.\\
  \item
    Der Data Service speichert bzw. ergänzt die Embeddings in
    \textbf{Milvus} an der passenden Stelle (Verknüpfung zu
    \texttt{document\_id}, \texttt{chunk\_id}, \texttt{tenant\_id}).\\
  \item
    In Postgres wird der Status des Dokuments (oder des gesamten
    Batches) auf \texttt{INDEXED} bzw. \texttt{COMPLETED} gesetzt.
  \end{itemize}
\item
  \textbf{Retrieval / Abfragephase}

  \begin{itemize}
  \tightlist
  \item
    Sobald der Batch fertig indexiert ist, kann der Client über die
    Middleware \textbf{Fragen} stellen.\\
  \item
    Die Inference-Pipeline aus Abschnitt 4.2 und das Diagramm 3.2
    greifen auf die in Milvus gespeicherten Embeddings und Chunks zu und
    liefern kontextualisierte Antworten zurück.
  \end{itemize}
\end{enumerate}

\subsubsection{\texorpdfstring{5.5 Mandantentrennung \& Deduplizierung
(\texttt{tenantId} \&
\texttt{contentHash})}{5.5 Mandantentrennung \& Deduplizierung (tenantId \& contentHash)}}\label{mandantentrennung-deduplizierung-tenantid-contenthash}

\textbf{Kurz-Zusammenfassung:}\\
Dieser Abschnitt beschreibt, wie \textbf{Mandantentrennung} über
\texttt{tenantId} und \textbf{Deduplizierung} über \texttt{contentHash}
in der Pipeline technisch umgesetzt werden.

\textbf{Technische Details}

\begin{itemize}
\tightlist
\item
  \textbf{Kernidee:}

  \begin{itemize}
  \tightlist
  \item
    \texttt{tenantId} bestimmt, \textbf{in welchen
    Partitionen/Collections} Daten gespeichert und später abgefragt
    werden.\\
  \item
    \texttt{contentHash} (z.\,B. \texttt{SHA-256(binary)}) identifiziert
    \textbf{identische Dokumentinhalte} innerhalb eines Tenants.
  \end{itemize}
\end{itemize}

\begin{figure}
\centering
\pandocbounded{\includesvg[keepaspectratio]{diagrams/diagram-09.svg}}
\caption{Diagramm 9}
\end{figure}

\textbf{Felder und Verwendung}

{\def\LTcaptype{none} % do not increment counter
\begin{longtable}[]{@{}
  >{\raggedright\arraybackslash}p{(\linewidth - 4\tabcolsep) * \real{0.1589}}
  >{\raggedright\arraybackslash}p{(\linewidth - 4\tabcolsep) * \real{0.1776}}
  >{\raggedright\arraybackslash}p{(\linewidth - 4\tabcolsep) * \real{0.6636}}@{}}
\toprule\noalign{}
\begin{minipage}[b]{\linewidth}\raggedright
Feld
\end{minipage} & \begin{minipage}[b]{\linewidth}\raggedright
Ebene
\end{minipage} & \begin{minipage}[b]{\linewidth}\raggedright
Zweck
\end{minipage} \\
\midrule\noalign{}
\endhead
\bottomrule\noalign{}
\endlastfoot
\textbf{\texttt{tenantId}} & Request, Topics, DB, Milvus & Strikte
Mandantentrennung über alle Stufen (S3-Pfad, Postgres,
Milvus-Collection/Partition). \\
\textbf{\texttt{contentHash}} & Data Service / Postgres & Erkennen von
Duplikaten je \texttt{tenantId}; Vermeiden doppelter Ingestion teurer
Pipelines. \\
\textbf{\texttt{documentId}} & Postgres, Topics, Milvus & Eindeutige
Identifikation je Dokument (auch bei mehreren logischen Referenzen). \\
\textbf{\texttt{batchId}} & Request, Postgres & Logische Gruppierung von
Uploads (z.\,B. Tages-/Monatsläufe eines Mandanten). \\
\end{longtable}
}

\textbf{Operational Note (Monitoring \& Fehlerszenarien)}

\begin{itemize}
\tightlist
\item
  \textbf{Monitoring-Signale}

  \begin{itemize}
  \tightlist
  \item
    \textbf{Deduplizierungsrate}: Verhältnis
    \texttt{duplicatesSkipped\ /\ acceptedCount} pro \texttt{tenantId}
    und \texttt{batchId}.\\
  \item
    \textbf{Anteil neuer \texttt{contentHash}-Werte} pro Zeitraum: Hoher
    Anteil kann auf \textbf{neue Dokumenttypen} oder fehlerhafte
    Hash-Berechnung hinweisen.
  \end{itemize}
\item
  \textbf{Typische Fehlerszenarien}

  \begin{itemize}
  \tightlist
  \item
    \textbf{Falsche \texttt{tenantId} im Upload} → Dokument landet in
    falscher Partition / falschem Index; Risiko von Daten-Leaks.\\
  \item
    \textbf{Fehlerhafte Hash-Berechnung} (z.\,B. unterschiedliche
    Normalisierung von Binärdaten) → Duplikate werden nicht erkannt,
    Ingestion-Kosten steigen.\\
  \item
    \textbf{Kollisionsrisiko} ist bei SHA-256 praktisch
    vernachlässigbar, sollte aber in
    Sicherheits-/Compliance-Diskussionen dokumentiert sein.
  \end{itemize}
\end{itemize}

\begin{center}\rule{0.5\linewidth}{0.5pt}\end{center}

\subsection{6. Warum dieses Konzept?
(Argumentation)}\label{warum-dieses-konzept-argumentation}

\begin{itemize}
\tightlist
\item
  \textbf{Entkopplung:} Wenn der Miner-Service ein Update erhält und
  kurz offline ist, füllen sich die Kafka-Topics, aber kein Dokument
  geht verloren.
\item
  \textbf{Kostenkontrolle (Azure):} Durch die Pufferung in Kafka
  verhindern wir, dass zu viele gleichzeitige Anfragen das Rate-Limit
  von Azure OpenAI sprengen.
\item
  \textbf{Audit-Log:} Jedes Kafka-Topic dient als permanentes Log. Wir
  können jederzeit nachvollziehen, wann ein Dokument welche
  Verarbeitungsstufe erreicht hat.
\item
  \textbf{Proxy-Performance:} Die Middleware (Spring Boot) bleibt extrem
  leichtgewichtig, da die rechenintensive Arbeit (KI \& Extraktion) in
  die Python-Worker ausgelagert ist.
\end{itemize}

\textbf{Takeaways:}

\begin{itemize}
\tightlist
\item
  \textbf{Warum Kafka + Microservices?} Entkopplung, Robustheit und
  Nachvollziehbarkeit.
\item
  \textbf{Warum Azure OpenAI + Milvus?} Trennung von LLM-Inferenz und
  Vektorspeicher mit klaren Verantwortlichkeiten.
\item
  \textbf{Warum mandantenfähig?} Saubere Isolation pro \texttt{tenantId}
  in allen Topics, Datenbanken und Indexen.
\end{itemize}

\textbf{Industrie-Kontext (RAG + Event-Streaming):}\\
Eine event-getriebene RAG-Architektur mit Kafka entspricht dem
etablierten Muster, \textbf{LLMs mit echtzeitnahem, kontextspezifischem
Wissen zu versorgen} und so Halluzinationen zu reduzieren. Kafka
übernimmt dabei die Rolle des \textbf{Datenrückgrats} für Ingestion,
Verarbeitung und konsistente Anreicherung -- vergleichbar mit dem Ansatz
„Real-Time GenAI with RAG using Apache Kafka and Flink`` (siehe
\hyperref[9-weiterfuxfchrende-literatur--referenzen]{Weiterführende
Literatur}).

\begin{center}\rule{0.5\linewidth}{0.5pt}\end{center}

\subsection{7. Betrieb, Skalierung \&
Bottlenecks}\label{betrieb-skalierung-bottlenecks}

\subsubsection{7.1 Wie wir Topics
verstehen}\label{wie-wir-topics-verstehen}

\begin{itemize}
\tightlist
\item
  \textbf{Append-only:} Nachrichten in Kafka-Topics werden \textbf{nicht
  editiert}, sondern nur angehängt.\\
\item
  \textbf{Abarbeitung:} Ein Dokument gilt in einem Verarbeitungsschritt
  als „abgearbeitet``, wenn der zuständige Service die Nachricht gelesen
  und seinen \textbf{Offset} committet hat.\\
\item
  \textbf{Historie:} Über die Retention bleiben Events für Analyse,
  Debugging und Replays verfügbar.
\end{itemize}

Damit bildet jedes Topic eine \textbf{Verarbeitungsstufe} der Pipeline
ab (z.\,B. „extrahiert``, „angereichert``, „vektorisiert``).

\subsubsection{7.2 Ablauf bei Multi-Upload (viele Dokumente auf
einmal)}\label{ablauf-bei-multi-upload-viele-dokumente-auf-einmal}

\begin{enumerate}
\def\labelenumi{\arabic{enumi}.}
\tightlist
\item
  Der Nutzer lädt \textbf{mehrere Dokumente} über die Middleware hoch.\\
\item
  Die Middleware erzeugt \textbf{pro Dokument} ein Event in
  \texttt{document-received} (optional mit gemeinsamer
  \texttt{batch\_id}).\\
\item
  Der Data Service liest diese Events, speichert die Rohdaten in S3 und
  erzeugt \textbf{pro Dokument} ein Event in
  \texttt{document-to-extract}.\\
\item
  Jeder weitere Schritt (Extraktion, Mining, Embeddings, Index) erzeugt
  wiederum \textbf{ein eigenes Event pro Dokument} in den jeweiligen
  Topics (\texttt{content-extracted}, \texttt{metadata-enriched},
  \texttt{vector-ready-to-index}).
\end{enumerate}

Durch diese feine Granularität können wir:

\begin{itemize}
\tightlist
\item
  bei Rückständen gezielt \textbf{eine Stage} skalieren (z.\,B. nur
  Extraction),
\item
  einzelne Dokumente oder Batches \textbf{nachverfolgen und
  reprocessen}.
\end{itemize}

\subsubsection{7.3 Retry-Mechanismen \&
Fehlerbehandlung}\label{retry-mechanismen-fehlerbehandlung}

Um Robustheit zu gewährleisten, arbeiten alle Services mit einem
einheitlichen Retry-Konzept:

\begin{itemize}
\tightlist
\item
  \textbf{Technische Fehler (z.\,B. Netzwerk, 5xx von Azure, temporäre
  S3-Probleme)}

  \begin{itemize}
  \tightlist
  \item
    Jeder Service führt \textbf{mehrere Retries mit Exponential Backoff}
    durch (z.\,B. 3--5 Versuche mit wachsender Wartezeit).\\
  \item
    Während der Retries wird \textbf{kein Offset committet}, d.\,h. die
    Message bleibt im Topic „offen``.\\
  \item
    Nach erfolgreichem Versuch wird der Offset einmalig committet.
  \end{itemize}
\item
  \textbf{Fachliche Fehler (z.\,B. defektes PDF, ungültiges Format,
  dauerhaftes 4xx)}

  \begin{itemize}
  \tightlist
  \item
    Nach einem konfigurierbaren Limit (\texttt{maxAttempts}) wird die
    Message in ein \textbf{Error-/DLQ-Topic} verschoben, z.\,B.:

    \begin{itemize}
    \tightlist
    \item
      \texttt{document-to-extract.dlq}\strut \\
    \item
      \texttt{content-extracted.dlq}\strut \\
    \item
      \texttt{vector-ready-to-index.dlq}\strut \\
    \end{itemize}
  \item
    Im DLQ-Topic wird zusätzlich ein \texttt{errorCode},
    \texttt{errorMessage} und ein \texttt{lastService}-Feld gespeichert.
  \end{itemize}
\item
  \textbf{Idempotenz}

  \begin{itemize}
  \tightlist
  \item
    Alle Operationen sind so ausgelegt, dass sie \textbf{idempotent}
    sind (z.\,B. „Embedding für \texttt{documentId\ +\ chunkId}
    existiert bereits`` → Skip / Update).\\
  \item
    Dadurch sind erneute Verarbeitungsversuche (Replays) jederzeit
    gefahrlos möglich -- sowohl automatisiert als auch manuell aus
    DLQ-Topics.
  \end{itemize}
\item
  \textbf{Monitoring \& Operations}

  \begin{itemize}
  \tightlist
  \item
    DLQ-Topics werden per Dashboard überwacht (Anzahl Messages,
    Fehlercodes).\\
  \item
    Über Admin-Tools können Einträge aus DLQs \textbf{manuell
    analysiert, korrigiert oder erneut eingespielt} werden (z.\,B. nach
    Fix eines Bugs oder einer Schema-Änderung).
  \end{itemize}
\end{itemize}

\subsubsection{7.4 Typische Bottlenecks pro
Stage}\label{typische-bottlenecks-pro-stage}

\begin{itemize}
\tightlist
\item
  \textbf{Extraction Service} (\texttt{document-to-extract} →
  \texttt{content-extracted})

  \begin{itemize}
  \tightlist
  \item
    CPU/GPU-bound (PDF-Parsing, OCR).\\
  \item
    Bottleneck-Signal: steigender Lag im Topic
    \texttt{document-to-extract}, hohe Auslastung der Extraction-Pods.\\
  \item
    Maßnahme: mehr Replikas des Extraction Services, Worker-Prozesse
    begrenzen.
  \end{itemize}
\item
  \textbf{Miner Service} (\texttt{content-extracted} →
  \texttt{metadata-enriched})

  \begin{itemize}
  \tightlist
  \item
    CPU-bound (Regex, NLP, fachliche Regeln).\\
  \item
    Bottleneck-Signal: Lag in \texttt{content-extracted}.\\
  \item
    Maßnahme: Scale-out, ggf. Regeln optimieren/batchen.
  \end{itemize}
\item
  \textbf{AI Service / Azure OpenAI} (\texttt{metadata-enriched} →
  \texttt{vector-ready-to-index})

  \begin{itemize}
  \tightlist
  \item
    Bottleneck: \textbf{Azure-Rate-Limits} (Requests/Minute,
    Tokens/Minute).\\
  \item
    Bottleneck-Signal: viele 429/5xx von Azure, wachsender Lag in
    \texttt{metadata-enriched}.\\
  \item
    Maßnahmen:

    \begin{itemize}
    \tightlist
    \item
      internes \textbf{Rate-Limiting} pro Pod (Token-Bucket, Queue),
    \item
      Concurrency pro Pod begrenzen,
    \item
      Scale-out nur bis zu einer Obergrenze, die zu den Azure-Limits
      passt.
    \end{itemize}
  \end{itemize}
\item
  \textbf{Milvus / Indexing} (\texttt{vector-ready-to-index})

  \begin{itemize}
  \tightlist
  \item
    Bottleneck: viele Inserts / Index-Operationen.\\
  \item
    Bottleneck-Signal: steigende Latenzen beim Schreiben, Lag in
    \texttt{vector-ready-to-index}.\\
  \item
    Maßnahmen: Batching von Vektoren, eigener Indexer-Service mit
    eigener Skalierung.
  \end{itemize}
\end{itemize}

\subsubsection{7.5 Zusammenspiel mit
Kubernetes}\label{zusammenspiel-mit-kubernetes}

\begin{itemize}
\tightlist
\item
  \textbf{Ein Deployment pro Service} (Middleware, Data, Extraction,
  Miner, AI, Indexer).\\
\item
  Alle Instanzen eines Services gehören in \textbf{eine Consumer-Gruppe}
  pro Topic -- zusätzliche Pods bedeuten mehr parallele Consumer.\\
\item
  \textbf{Auto-Scaling nach Kafka-Lag:}

  \begin{itemize}
  \tightlist
  \item
    Z.\,B. via KEDA/HPA, die auf „Lag pro Consumer-Gruppe`` schauen
    (z.\,B.
    \texttt{lag(document-to-extract,\ extraction-service-group)}).\\
  \item
    So werden genau die Stages hochskaliert, in denen sich Messages
    stauen.
  \end{itemize}
\item
  \textbf{Externe Limits beachten:}

  \begin{itemize}
  \tightlist
  \item
    Beim AI Service verhindern interne Limits, dass
    Kubernetes-Skalierung die \textbf{Azure-Rate-Limits} überfährt.\\
  \item
    Durch Idempotenz (z.\,B. pro \texttt{doc\_id}+\texttt{chunk\_id})
    bleiben Replays und Re-Processing robust.
  \end{itemize}
\end{itemize}

Damit ist klar dokumentiert, \textbf{wie} unsere Topics genutzt werden,
\textbf{wo} Bottlenecks entstehen können und \textbf{wie} wir sie über
Kubernetes und Kafka-Lag im Betrieb steuern.

\subsubsection{7.6 Fehlerszenarien \&
Reprocessing-Strategien}\label{fehlerszenarien-reprocessing-strategien}

\textbf{Kurz-Zusammenfassung:}\\
Dieser Abschnitt fasst typische Fehlerszenarien in der RAG-Pipeline
zusammen und beschreibt, wie über \textbf{Retries}, \textbf{DLQs} und
\textbf{Reprocessing} stabiler Betrieb sichergestellt wird.

\textbf{Typische Fehlerszenarien (Auszug)}

{\def\LTcaptype{none} % do not increment counter
\begin{longtable}[]{@{}
  >{\raggedright\arraybackslash}p{(\linewidth - 6\tabcolsep) * \real{0.2000}}
  >{\raggedright\arraybackslash}p{(\linewidth - 6\tabcolsep) * \real{0.2579}}
  >{\raggedright\arraybackslash}p{(\linewidth - 6\tabcolsep) * \real{0.2947}}
  >{\raggedright\arraybackslash}p{(\linewidth - 6\tabcolsep) * \real{0.2474}}@{}}
\toprule\noalign{}
\begin{minipage}[b]{\linewidth}\raggedright
Szenario
\end{minipage} & \begin{minipage}[b]{\linewidth}\raggedright
Ursache (typisch)
\end{minipage} & \begin{minipage}[b]{\linewidth}\raggedright
Signal im Betrieb
\end{minipage} & \begin{minipage}[b]{\linewidth}\raggedright
Maßnahme / Strategie
\end{minipage} \\
\midrule\noalign{}
\endhead
\bottomrule\noalign{}
\endlastfoot
\textbf{Azure 5xx/429 beim Embedding} & Azure-Störung, Rate-Limits
überschritten & Viele 5xx/429 im AI-Service, wachsender Lag in
\texttt{metadata-enriched} & Retries mit Backoff, internes
Rate-Limiting, ggf. Pods drosseln statt skalieren. \\
\textbf{Defektes PDF / inkompatibles Format} & Kaputte Datei, nicht
unterstütztes Format & Wiederholte Fehler im Extraction Service,
Einträge im \texttt{document-to-extract.dlq} & Nach X Retries in DLQ
verschieben, Datei/Quelle korrigieren, gezieltes Reprocessing. \\
\textbf{Fachliche Parsing-Fehler} & Neue Layouts/Belege, nicht
abgedeckte Regeln & Anstieg von \texttt{errorCode} in Miner-DLQs,
Fach-Feedback & Regeln/Modelle anpassen, DLQ-Einträge nach Fix erneut
einspielen. \\
\textbf{Milvus-Write-Latenzen} & Hohe Schreiblast, fehlendes Batching &
Latenzspitzen beim Indexing, Lag in \texttt{vector-ready-to-index} &
Batching erhöhen, eigenen Indexer-Service skalieren. \\
\textbf{Fehlerhafte Mandanten-Konfiguration} & Falsche
\texttt{tenantId}/Partition & Dokumente fehlen in erwarteten Antworten,
Inkonsistenzen in Audits & Korrektur in Postgres, ggf. Re-Indexing
betroffener Dokumente. \\
\end{longtable}
}

\textbf{Reprocessing-Flows (vereinfacht)}

\begin{figure}
\centering
\pandocbounded{\includesvg[keepaspectratio]{diagrams/diagram-10.svg}}
\caption{Diagramm 10}
\end{figure}

\textbf{Operational Note}

\begin{itemize}
\tightlist
\item
  \textbf{Monitoring-Perspektive (SRE/DevOps)}

  \begin{itemize}
  \tightlist
  \item
    DLQ-Größen und -Wachstumsgeschwindigkeit pro Topic und
    \texttt{tenantId}.\\
  \item
    Kafka-Lag pro Consumer-Gruppe als Indikator für \textbf{Staus} in
    einzelnen Stages.\\
  \item
    Fehlerraten nach Service (Extraction, Miner, AI, Indexer) und nach
    Fehlertyp (\texttt{errorCode}).
  \end{itemize}
\item
  \textbf{Best Practices}

  \begin{itemize}
  \tightlist
  \item
    Reprocessing \textbf{immer idempotent} gestalten (z.\,B. Embeddings
    pro \texttt{documentId\ +\ chunkId} überschreibbar).\\
  \item
    Replays \textbf{batchweise} durchführen und Last auf Azure/Milvus
    berücksichtigen (Zeitfenster, Rate-Limits).\\
  \item
    Bei kritischen Tenants ggf. \textbf{separates Maintenance-Fenster}
    für großvolumige Re-Indexing-Jobs vorsehen.
  \end{itemize}
\end{itemize}

\begin{center}\rule{0.5\linewidth}{0.5pt}\end{center}

\subsection{8. Glossar (wichtige
Begriffe)}\label{glossar-wichtige-begriffe}

\begin{itemize}
\tightlist
\item
  \textbf{Middleware}: Spring-Boot-Service als API-Gateway/Proxy; nimmt
  Requests vom Client entgegen, kümmert sich um Authentifizierung und
  reicht Aufrufe an die internen Services weiter.
\item
  \textbf{Data Service}: Python-Service zur Orchestrierung der
  Datenflüsse (S3, Postgres, Milvus) und Verwaltung von
  Dokument‑/Batch-Status.
\item
  \textbf{Extraction Service}: Python-Service für die Extraktion von
  Text und Struktur aus binären Dokumenten (PDF, Bilder, Office) via
  Azure Document Intelligence.
\item
  \textbf{Miner Service}: Python-Service zur fachlichen Anreicherung
  (z.\,B. Erkennen von Rechnungsnummern, Beträgen, IDs) auf Basis des
  extrahierten Textes.
\item
  \textbf{AI Service}: Python-Service für Chunking, Embedding-Erzeugung
  (Azure OpenAI) und LLM-Aufrufe für Antworten.
\item
  \textbf{Milvus}: Vektordatenbank für Speicherung und Ähnlichkeitssuche
  der Embeddings; bildet die Grundlage der Wissenssuche.
\item
  \textbf{Kafka Topic}: Append-only Log, in das Services Events
  schreiben (Producer) und aus dem andere Services lesen (Consumer);
  bildet jeweils eine Verarbeitungsstufe der Pipeline ab.
\item
  \textbf{DLQ (Dead Letter Queue)}: Spezielles Kafka-Topic, in das
  Nachrichten mit dauerhaften Fehlern verschoben werden; dient zur
  Analyse und zum manuellen/automatisierten Re-Processing.
\item
  \textbf{\texttt{tenantId}}: Technische Mandanten-ID; wird in allen
  Topics, Datenbanken und Indexen mitgeführt, um Daten strikt
  mandantengetrennt zu halten.
\item
  \textbf{\texttt{documentId}}: Eindeutige ID eines Dokuments innerhalb
  eines Tenants; verknüpft Rohdatei, Extraktion, Metadaten, Chunks und
  Embeddings.
\item
  \textbf{\texttt{batchId}}: Optionale ID, die mehrere Dokumente eines
  Bulk-Uploads logisch gruppiert (z.\,B. ein Tages- oder Monatslauf
  eines Mandanten).
\item
  \textbf{\texttt{contentHash}}: Hash-Wert (z.\,B. SHA‑256) über den
  Dateiinhalt; dient zur Erkennung von Duplikaten über mehrere Uploads
  hinweg.
\item
  \textbf{Chunk}: Kleinere Texteinheit eines Dokuments (Abschnitt,
  Absatz etc.), die einzeln eingebettet und in Milvus gespeichert wird.
\item
  \textbf{Embedding}: Hochdimensionaler Vektor, der die Bedeutung eines
  Chunks oder einer Frage für die Vektorsuche repräsentiert.
\end{itemize}

\begin{center}\rule{0.5\linewidth}{0.5pt}\end{center}

\subsection{9. Weiterführende Literatur \&
Referenzen}\label{weiterfuxfchrende-literatur-referenzen}

{\def\LTcaptype{none} % do not increment counter
\begin{longtable}[]{@{}
  >{\raggedright\arraybackslash}p{(\linewidth - 2\tabcolsep) * \real{0.2963}}
  >{\raggedright\arraybackslash}p{(\linewidth - 2\tabcolsep) * \real{0.7037}}@{}}
\toprule\noalign{}
\begin{minipage}[b]{\linewidth}\raggedright
Quelle
\end{minipage} & \begin{minipage}[b]{\linewidth}\raggedright
Inhalt / Relevanz
\end{minipage} \\
\midrule\noalign{}
\endhead
\bottomrule\noalign{}
\endlastfoot
\textbf{\href{https://www.kai-waehner.de/blog/2024/05/30/real-time-genai-with-rag-using-apache-kafka-and-flink-to-prevent-hallucinations/}{Real-Time
GenAI with RAG using Apache Kafka and Flink to Prevent Hallucinations}}
(Kai Waehner, Mai 2024) & Erläutert, warum \textbf{RAG + Event-Streaming
(Kafka/Flink)} LLMs mit echtzeitnahem, kontextspezifischem Datenzugriff
versorgt und so \textbf{Halluzinationen reduziert}. Themen: Real-Time
Data Ingestion, dynamisches Indexing, Skalierbarkeit, kontextuelle
Anreicherung in der RAG-Pipeline; inkl. Lightboard-Video und
\textbf{Expedia-Case-Study} (conversational chatbot, 60 \%+
Self-Service, 40 \%+ Kosteneinsparung bei Agenten). Direkt
anschlussfähig an unsere Architektur (Kafka als Event-Bus,
Mandantenfähigkeit, Ingestion vs.~Inference). \\
\end{longtable}
}
